

\section{La Lecture Extensive : Outil roi de la consolidation}
space{0.3cm}
oindent

L'enseignement du grec ancien se trouve aujourd'hui à la croisée des chemins, tiraillé entre une
tradition philologique séculaire, axée sur le décodage analytique (Grammaire-Traduction), et les
apports révolutionnaires de la psycholinguistique et des sciences de l'acquisition des langues
secondes (SLA). Le présent rapport de recherche se concentre sur le point III.3.\cite{III-3-2-ref2}
de la thèse, consacré à la \textbf{Lecture Extensive (Extensive Reading ou ER)}. Cette approche,
loin d'être une simple méthode complémentaire, constitue une refonte fondamentale de la manière dont
le cerveau apprend à traiter une langue morte. Elle postule que l'automatisation des processus de
reconnaissance linguistique — condition \textit{sine qua non} d'une lecture fluide et sans effort —
ne peut s'acquérir par l'étude consciente des règles, mais nécessite une exposition massive à des
textes compréhensibles, spécifiquement conçus pour respecter les contraintes cognitives de
l'apprenant : les \textit{Graded Readers} (lectures graduées).

L'enjeu n'est pas trivial. La méthode traditionnelle produit souvent des étudiants capables de
disserter sur la structure de l'aoriste mais incapables de lire une page de Platon sans s'arrêter à
chaque ligne pour consulter un dictionnaire ou un tableau de paradigmes. Ce déficit de fluidité
n'est pas un échec intellectuel, mais le résultat d'une surcharge cognitive systémique. En analysant
les mécanismes neurocognitifs de l'automatisation (LaBerge \& Samuels, Ullman), les statistiques de
la couverture lexicale (Nation), et l'état de l'art bibliographique des ressources disponibles (de
Rouse à GlossaHouse), on peut démontrer que la lecture extensive est la technologie pédagogique
manquante pour transformer le déchiffrage laborieux en lecture véritable.



\subsection{1\. fondements neurocognitifs de l'automatisation de la lecture}

Pour comprendre pourquoi la lecture extensive est indispensable, il faut d'abord disséquer ce qui se
passe dans le cerveau d'un lecteur fluide par opposition à un lecteur débutant ou \og
déchiffreur\fg{}. La littérature scientifique suggère que la fluidité n'est pas une accélération du
processus de traduction consciente, mais le basculement vers un système de traitement de
l'information radicalement différent.



\subsection{La théorie de l'automaticité de laberge et samuels}

Le cadre théorique dominant pour expliquer la fluidité en lecture est la théorie de l'automatisation
du traitement de l'information, formulée par LaBerge et Samuels (1974). Ce modèle part d'un postulat
économique : l'esprit humain dispose d'une quantité limitée d'énergie
attentionnelle.\cite{III-3-2-ref1} La lecture est une tâche complexe qui implique plusieurs sous-
processus simultanés : le décodage visuel des lettres, la reconnaissance des mots, l'analyse
syntaxique et l'intégration sémantique.

Si le lecteur doit allouer son attention consciente aux processus de bas niveau (décodage des
lettres, identification des formes verbales), il sature sa mémoire de travail. Il ne reste alors
plus de ressources cognitives pour les processus de haut niveau : la compréhension du sens,
l'inférence contextuelle et l'appréciation esthétique.\cite{III-3-2-ref1} C'est la situation typique
de l'étudiant en grec ancien formé par la méthode Grammaire-Traduction : face à une forme comme
\textit{ἐπαιδευσάμην}, il mobilise son attention pour identifier l'augment, la racine, le suffixe
sigmatique et la désinence moyenne. Ce traitement est sériel, lent et épuisant.

L'automatisation, visée par la lecture extensive, se définit comme la capacité à exécuter ces tâches
de bas niveau sans attention volontaire. Un lecteur fluide reconnaît le mot \textit{ἄνθρωπος} aussi
instantanément qu'il reconnaît un visage familier, sans passer par une analyse des composants. Selon
Logan (1997), cette automatisation est caractérisée par la vitesse, l'absence d'effort et
l'autonomie du processus.\cite{III-3-2-ref1} LaBerge et Samuels ont démontré que la répétition
massive est le mécanisme clé de cette transition : c'est en rencontrant les mêmes motifs visuels et
morphologiques des milliers de fois que le cerveau construit des représentations unitaires
accessibles directement, libérant ainsi l'attention pour le sens.\cite{III-3-2-ref2}



\subsection{Le modèle déclaratif/procédural (dp) et le cerveau bilingue}

Les travaux de Michael Ullman sur le modèle Déclaratif/Procédural (DP) offrent une explication
neurobiologique à cette distinction entre \og savoir grammatical\fg{} et \og savoir lire\fg{}.
Ullman postule que le langage repose sur deux systèmes de mémoire à long terme distincts, ancrés
dans des substrats neuronaux différents :

\begin{table}[h!]\small\centering\begin{tabular}{| p{2cm} | p{2.2cm} | p{2.2cm} | p{2.2cm} |}\hline\textbf{Système de Mémoire} & \textbf{Substrat Neuronal} & \textbf{Fonction Linguistique} & \textbf{Caractéristiques} \\ \hline\textbf{Mémoire Déclarative} & Lobe temporal médial (Hippocampe) & Lexique mental, faits, règles explicites & Apprentissage rapide, conscient, mais traitement lent. \\ \hline\textbf{Mémoire Procédurale} & Ganglions de la base (Striatum), Aire de Broca & Grammaire, syntaxe, séquences motrices & Apprentissage lent (par répétition), inconscient, traitement ultra-rapide. \\ \hline\end{tabular}\end{table}Dans l'apprentissage traditionnel du grec (L2), les étudiants stockent la grammaire dans leur
mémoire déclarative (apprendre \og par cœur\fg{} des tableaux). Or, la mémoire déclarative n'est pas
optimisée pour le traitement en temps réel requis par la lecture fluide ou la
parole.\cite{III-3-2-ref3} La lecture extensive fonctionne comme un \og entraînement
procédural\fg{}. En exposant le cerveau à des milliers d'instanciations correctes de structures
grammaticales (par exemple, l'usage de l'accusatif pour l'objet direct), l'apprenant transfère
progressivement le traitement de ces structures du système déclaratif vers le système
procédural.\cite{III-3-2-ref5}

Ce transfert est crucial car le système procédural est le seul capable de gérer la complexité
combinatoire du grec ancien à la vitesse de la lecture naturelle. Ullman et ses collaborateurs
soulignent que si l'enseignement explicite peut initier le processus (déclaratif), seule la pratique
massive (input) peut consolider les circuits procéduraux.\cite{III-3-2-ref6} Sans lecture extensive,
l'étudiant reste piégé dans un mode de traitement déclaratif, condamné à une lecture lente et
cognitivement coûteuse.



\subsection{La double voie et la décomposition morphologique}

La spécificité du grec ancien, langue hautement flexionnelle et fusionnelle, pose des défis
particuliers au modèle de lecture. Contrairement à l'anglais où la morphologie est pauvre, le grec
encode une densité d'information extrême dans chaque mot. Le \og Dual Route Model\fg{} (Modèle à
Double Voie) de la lecture suggère deux chemins d'accès au sens :

1. \textbf{La Voie Directe (Whole-word)} : Reconnaissance globale du mot comme une image
orthographique.

2. \textbf{La Voie de Décomposition} : Analyse des constituants (racine \+
affixes).\cite{III-3-2-ref8}

Pour les lecteurs natifs de langues riches morphologiquement, la voie de décomposition est très
efficace et automatisée. Cependant, pour l'apprenant L2, la décomposition est souvent un goulot
d'étranglement. Des études utilisant la magnétoencéphalographie (MEG) et des tâches d'amorçage ont
montré que les apprenants L2 traitent souvent les formes fléchies complexes comme des mots entiers
(stockage lexical) faute de pouvoir décomposer automatiquement, ou inversement, s'épuisent dans une
décomposition analytique lente.\cite{III-3-2-ref10}

La lecture extensive de textes simplifiés (\textit{Graded Readers}) permet d'entraîner
spécifiquement cette décomposition morphologique rapide. En voyant le suffixe \textit{\-o-men}
attaché à des centaines de verbes différents dans des contextes clairs, le cerveau apprend à traiter
ce suffixe comme une unité de sens distincte mais automatiquement intégrée, réduisant ainsi ce que
l'on appelle la \og charge fovéale\fg{}.\cite{III-3-2-ref12} Une charge fovéale trop élevée (un mot
trop complexe à traiter) empêche le lecteur d'utiliser sa vision parafovéale pour anticiper le mot
suivant, ce qui brise le flux de la lecture.\cite{III-3-2-ref14}



\subsection{L'importance de la prosodie}

Enfin, il convient de noter une critique importante apportée au modèle de l'automaticité pure :
celle de Schreiber (1980, 1987). L'automatisation de la reconnaissance des mots ne suffit pas si
elle n'est pas accompagnée d'une structuration prosodique (rythme, intonation, pause). La prosodie
est l'interface entre le décodage lexical et la syntaxe.\cite{III-3-2-ref16} En grec ancien, cela
implique que la lecture extensive gagne à être accompagnée d'une oralisation (ou d'une
subvocalisation correcte) et d'écoute (audiobooks), pour que la \og musique\fg{} de la phrase aide à
segmenter le flux d'informations et à maintenir le sens en mémoire de travail.\cite{III-3-2-ref17}



\subsection{2\. les métriques de la compréhension : combien de mots faut-il?}

Si l'automatisation nécessite une lecture massive sans effort, comment définir \og sans effort\fg{}?
Les recherches en linguistique appliquée, notamment celles de Paul Nation, ont quantifié avec
précision les seuils de vocabulaire nécessaires.



\subsection{La règle impérative des 98\%}

L'étude fondamentale de Hu et Nation (2000) a établi qu'un lecteur doit connaître \textbf{98\%} des
mots d'un texte pour pouvoir le lire confortablement, sans aide extérieure, et en comprendre le sens
global tout en devinant les mots inconnus par le contexte.\cite{III-3-2-ref18}

\begin{itemize}\item \textbf{À 98\% de couverture} (1 mot inconnu tous les 50 mots) : La lecture est fluide, le plaisir est possible, et l'apprentissage incident (incidental learning) du vocabulaire nouveau se produit naturellement.\item \textbf{À 95\% de couverture} (1 mot inconnu tous les 20 mots) : La compréhension baisse significativement. Le lecteur doit recourir à des stratégies de compensation conscientes. C'est le seuil minimal pour une lecture assistée intensive, mais insuffisant pour l'ER.\cite{III-3-2-ref19}\item \textbf{À 80-90\% de couverture} : C'est le niveau typique d'un étudiant de grec face à un texte authentique après un an d'étude. À ce niveau, la lecture est impossible; c'est un décryptage cryptographique.\cite{III-3-2-ref20}\end{itemize}Pour l'anglais, atteindre ce seuil de 98\% sur des romans non simplifiés requiert un vocabulaire de
8 000 à 9 000 familles de mots.\cite{III-3-2-ref20} Pour le grec ancien, la situation est encore
plus critique en raison de la nature close du corpus et de la distribution lexicale.



\subsection{Le paradoxe 98-56 du corpus grec}

Seumas Macdonald a mis en lumière une statistique effrayante pour les apprenants du Nouveau
Testament grec, qu'il nomme le \og Paradoxe 98-56\fg{}.\cite{III-3-2-ref23} Pour lire le Nouveau
Testament avec une couverture de 98\% (le seuil de fluidité), un étudiant doit connaître environ
\textbf{56\%} de tout le vocabulaire unique du corpus, soit plus de 3 000 mots. Or, les manuels
d'introduction typiques enseignent entre 300 et 500 mots, ce qui laisse l'étudiant à des années-
lumière de la lecture fluide du texte cible.

Ce fossé (gap) immense entre la fin du manuel et le début du texte authentique est la zone de la
mort pour la plupart des apprenants. Sans \textit{Graded Readers} pour combler cet espace — en
fournissant des textes à 98\% de couverture pour des vocabulaires de 500, 1000, 1500, 2000 mots —
l'automatisation ne peut jamais se produire, car l'étudiant est constamment en surcharge cognitive.



\subsection{Recyclage lexical et répétition espacée}

L'acquisition d'un mot n'est pas instantanée. La recherche indique qu'un mot doit être rencontré
entre 5 et 20 fois dans des contextes variés pour être intégré.\cite{III-3-2-ref24} Dans les textes
authentiques, de nombreux mots sont des \textit{hapax legomena} (apparaissent une seule fois) ou
sont très rares. La lecture de textes authentiques est donc inefficace pour l'acquisition initiale
du vocabulaire, car le taux de recyclage est trop faible.

Les \textit{Graded Readers} sont conçus pour manipuler artificiellement ce taux de répétition. En
limitant le vocabulaire (Sheltered Vocabulary), ils forcent la réapparition des mots cibles,
assurant les expositions multiples nécessaires à la mémorisation procédurale.\cite{III-3-2-ref25}
L'analyse des manuels traditionnels montre souvent un échec sur ce point, introduisant trop de
vocabulaire sans recyclage suffisant.\cite{III-3-2-ref26}



\subsection{3\. typologie et analyse bibliographique des ressources er}

La reconnaissance de ces principes a conduit à une renaissance de la production littéraire
pédagogique en grec ancien. Nous pouvons classer ces ressources en fonction de leur adhésion aux
principes de Day \& Bamford et de leur niveau de contrôle lexical.



\subsection{Les pionniers de la méthode directe (1900-1920)}

Bien avant la théorisation de la SLA, des pédagogues britanniques avaient intuitivement compris la
nécessité de textes simples.

\begin{itemize}\item \textbf{W.H.D. Rouse, \textit{A Greek Boy at Home} (1909)} : Écrit pour accompagner son enseignement oral à la Perse School, ce livre raconte la vie quotidienne d'un garçon athénien. Le vocabulaire est concret, répétitif et contextuel.\cite{III-3-2-ref27} C'est un ancêtre direct des novellas modernes.\item \textbf{F.H. Colson, \textit{Stories and Legends: A First Greek Reader} (1908)} : Offre des récits mythologiques simplifiés. Le contrôle lexical est moins rigoureux que dans les standards modernes, mais reste accessible.\cite{III-3-2-ref30}\item \textbf{C.M. Moss, \textit{A First Greek Reader} (1887)} : Fournit des anecdotes historiques et mythologiques avec un vocabulaire limité.\cite{III-3-2-ref31}\end{itemize}Ces ouvrages, souvent disponibles dans le domaine public, constituent une ressource précieuse mais
présentent parfois un vocabulaire vieilli ou une progression grammaticale trop abrupte pour
l'autodidacte moderne.\cite{III-3-2-ref29}



\subsection{Les manuels narratifs et la méthode "naturelle" (fin xxe)}

\begin{itemize}\item \textbf{Athenaze (Édition Italienne / Vivarium Novum)} : C'est la référence absolue actuelle pour l'immersion textuelle. Luigi Miraglia a considérablement étendu le texte narratif de l'édition originale (anglaise) d'Oxford. Le volume 1 italien contient environ 26 000 mots de texte grec continu (contre \~12 000 pour l'anglais), et avec les suppléments (\textit{Ephodion}, \textit{Meletemata}), l'étudiant est exposé à plus de 30 000 mots.\cite{III-3-2-ref27} La progression suit une histoire continue, permettant un attachement émotionnel et contextuel fort.\item \textbf{Alexandros : To Hellenikon Paidion} : Écrit par Mario Díaz Avila, cet ouvrage est conçu comme l'équivalent grec du célèbre \textit{Lingua Latina Per Se Illustrata} (LLPSI) de Hans Ørberg. Entièrement en grec, richement illustré, il permet une lecture sans traduction dès la première page. Bien que plus court que LLPSI (environ 5 000 mots), il offre une excellente introduction inductive.\cite{III-3-2-ref27}\item \textbf{JACT \textit{Reading Greek}} : Ce cours britannique repose sur une narration continue adaptée d'Aristophane et de Démosthène. Bien que la grammaire soit traitée séparément, le volume de texte est conséquent et conçu pour être lu pour le sens.\cite{III-3-2-ref28}\end{itemize}

\subsection{La révolution des novellas et du "sheltered vocabulary" (xxie siècle)}

Depuis 2015, une communauté d'enseignants-auteurs, influencée par les théories de Krashen et le
mouvement TPRS (Teaching Proficiency through Reading and Storytelling), produit des ouvrages
spécifiquement calibrés pour les principes de l'ER (couverture de 98\%, répétition massive, intérêt
narratif).



\subsection{A. les novellas et le mouvement "indie"}

Ces livres sont courts, ciblés sur un nombre très restreint de mots uniques (souvent moins de 200),
et visent la fluidité totale.

\begin{itemize}\item \textbf{Seumas Macdonald} : Auteur de \textit{Lingua Graeca Per Se Illustrata} (LGPSI) et de novellas comme \textit{A Spartan Tale}. Il applique rigoureusement le contrôle du vocabulaire pour permettre une lecture extensive dès le niveau débutant. Il défend une approche où le volume de lecture prime sur la complexité.\cite{III-3-2-ref33}\item \textbf{Lance Piantaggini (Magister P)} : Célèbre pour ses novellas latines (le \og Pisoverse\fg{}), il collabore et inspire la production grecque (ex: \textit{Markos Ho Magos}). Ses principes de design incluent l'usage massif de cognats, la limitation stricte des mots uniques (souvent \< 50 pour les premiers niveaux) et l'écriture de récits \og compelling\fg{} (captivants) qui ne sont pas de simples prétextes grammaticaux.\cite{III-3-2-ref35}\item \textbf{Juan Coderch} : Professeur à St Andrews, il a traduit des classiques modernes en grec attique : \textit{Le Petit Prince} (Ὁ Μικρὸς Πρίγκιψ), \textit{Sherlock Holmes}, \textit{Don Camillo}. Ses ouvrages innovent par leur mise en page : le vocabulaire difficile et les explications linguistiques sont fournis en marge, en grec simple, permettant de rester dans la langue cible sans recourir à la L1.\cite{III-3-2-ref33}\end{itemize}

\subsection{B. glossahouse et les ressources bibliques}

La maison d'édition \textbf{GlossaHouse} a industrialisé la production de matériel ER pour le grec
Koinè, comblant le fossé mentionné par Macdonald.

\begin{itemize}\item \textbf{AGROS Series} : Une série de ressources illustrées et graduées.\item \textbf{Tiered Readers} : Des ouvrages qui présentent le texte biblique ou patristique avec différents niveaux d'aide (vocabulaire glossé pour les mots apparaissant moins de 50 fois, puis 30 fois, etc.), permettant à l'étudiant de lire des textes authentiques avec un \og échafaudage\fg{} (scaffolding) adapté.\cite{III-3-2-ref39}\item \textbf{Ouvrages Illustrés} : Comme \textit{The GlossaHouse Illustrated Greek-English New Testament}, qui utilise l'image pour renforcer la mémorisation sémantique sans traduction.\cite{III-3-2-ref41}\end{itemize}

\subsection{Tableau synthétique des ressources er : typologie et niveaux}

Le tableau ci-dessous classe les ressources selon leur densité lexicale et leur fonction
pédagogique, permettant de construire un curriculum de lecture progressive.

\begin{table}[h!]\small\centering\begin{tabular}{| l | l | l | l | l |}\hline\textbf{Ressource / Auteur} & \textbf{Type} & \textbf{Dialecte} & \textbf{Niveau Cible} & \textbf{Caractéristique Méthodologique} \\ \hline\textbf{Alexandros} (Díaz Avila) & Manuel Illustré & Attique & Débutant (A1) & Méthode Naturelle, 100\% Grec, imitation LLPSI. \\ \hline\textbf{Athenaze} (Ed. Italienne) & Manuel Narratif & Attique & A1 \-\> B1 & Volume massif (\~30k mots), progression grammaticale inductive. \\ \hline\textbf{Seumas Macdonald} & Novellas (ex: \textit{Spartan Tale}) & Attique & A1 / A2 & Vocabulaire strictement contrôlé (\<200 mots uniques), répétition élevée. \\ \hline\textbf{GlossaHouse Readers} & Textes Gradués & Koinè & A2 \-\> C1 & Glosses basées sur la fréquence (ex: mots \<50 occurrences). \\ \hline\textbf{Juan Coderch} & Traductions Modernes & Attique & B1 / B2 & \textit{Le Petit Prince}, \textit{Sherlock Holmes}. Vocabulaire en marge, style idiomatique. \\ \hline\textbf{Geoffrey Steadman} & Textes Classiques Annotés & Attique / Homérique & B2 \-\> C1 & Texte authentique \+ Vocabulaire et Grammaire sur la même page. \\ \hline\textbf{W.H.D. Rouse} & Lecteur Historique & Attique & A2 & \textit{A Greek Boy at Home}. Immersion vie quotidienne, domaine public. \\ \hline\end{tabular}\end{table}

\subsection{4\. efficacité pédagogique et preuves empiriques}

L'engouement pour la lecture extensive en langues anciennes ne relève pas seulement d'une mode, mais
s'appuie sur des données empiriques probantes qui valident les théories de l'acquisition.



\subsection{Réduction de l'anxiété et motivation (étude humphreys 2022\)}

Une thèse doctorale récente de Dustin Humphreys (2022) a comparé quantitativement les résultats
affectifs et cognitifs d'étudiants en grec ancien suivant deux méthodologies : la méthode Grammaire-
Traduction (GTM) et l'Approche Communicative (CLT) incluant l'ER.

\begin{itemize}\item \textbf{Résultats} : L'étude a montré une différence statistiquement significative (p \< 0.05). Les étudiants du groupe CLT/ER affichaient une \textbf{auto-efficacité en lecture} (Reading Self-Efficacy) beaucoup plus élevée (Score moyen 8.\cite{III-3-2-ref21} vs 7.\cite{III-3-2-ref54} pour GTM) et une anxiété débilitante moindre.\cite{III-3-2-ref42}\item \textbf{Implication} : L'auto-efficacité est un prédicteur majeur de la réussite à long terme. Les étudiants qui se sentent capables de lire lisent davantage, créant un cercle vertueux d'exposition à l'input.\cite{III-3-2-ref42}\end{itemize}

\subsection{L'effet "clockwork orange" et l'acquisition incidentelle}

Une étude souvent citée pour valider l'ER est celle de l'acquisition du vocabulaire \og Nadsat\fg{}
(l'argot inventé par Anthony Burgess dans A Clockwork Orange). Les lecteurs acquièrent ce
vocabulaire simplement en lisant le roman, sans dictionnaire, grâce à la répétition contextuelle.
Les participants retenaient en moyenne 76\% du vocabulaire Nadsat après la
lecture.\cite{III-3-2-ref43}

Ce phénomène est reproduit avec les novellas grecques modernes : les étudiants acquièrent des mots
comme τράπεζα (table) ou βαδίζει (il marche) non pas en les apprenant par cœur, mais parce qu'ils
apparaissent 50 fois dans l'histoire de Alexandros ou d'une novella de Piantaggini. L'absence
d'effort conscient (learning vs acquisition) garantit que ces mots sont stockés de manière à être
récupérables instantanément.\cite{III-3-2-ref44}



\subsection{Les résultats institutionnels (polis \& vivarium novum)}

Les institutions qui ont adopté l'ER et l'immersion comme cœur de leur curriculum rapportent des
résultats supérieurs en termes de fluidité.

\begin{itemize}\item \textbf{Institut Polis (Jérusalem)} : En utilisant la méthode Polis (qui combine TPR et lecture extensive), les recherches internes montrent que les étudiants internalisent les structures grammaticales (comme les cas) beaucoup plus rapidement qu'avec une approche analytique. L'approche \og Listen First\fg{} prépare le cerveau à la lecture en établissant les représentations phonologiques nécessaires.\cite{III-3-2-ref45}\item \textbf{Vivarium Novum (Rome)} : Le curriculum exige la lecture de milliers de pages. Les étudiants, qui vivent en immersion latine et grecque, développent une capacité de lecture à vue (sight reading) qui dépasse largement les standards universitaires classiques, validant l'idée que le volume d'input est le facteur déterminant de la compétence.\cite{III-3-2-ref47}\end{itemize}

\subsection{La problématique de la prononciation et de l'oralisation}

Les études empiriques soulignent également l'importance de la composante orale pour soutenir la
lecture. Le \textbf{Biblical Language Center} a démontré que l'intégration de l'audio et de
l'interaction orale (communicative approach) améliore la rétention du vocabulaire à long terme par
rapport à la simple lecture silencieuse ou à la traduction.\cite{III-3-2-ref50} Cela confirme les
théories de Schreiber sur le rôle de la prosodie dans la fluidité de lecture : entendre la langue
aide à la lire.\cite{III-3-2-ref17}



\subsection{5\. stratégies de mise en œuvre : le fvr et l'échelonnage}

Comment implémenter concrètement l'ER dans un programme de grec ancien ou en autodidacte? Les
principes de Day \& Bamford doivent être adaptés à la réalité matérielle du corpus grec.



\subsection{La lecture libre volontaire (fvr) et la bibliothèque de classe}

Le principe de \og Free Voluntary Reading\fg{} (FVR) est central. L'enseignant doit constituer une
bibliothèque de classe contenant des exemplaires multiples des novellas mentionnées (Macdonald,
Piantaggini, Coderch, GlossaHouse).\cite{III-3-2-ref52}

\begin{itemize}\item \textbf{Sustained Silent Reading (SSR)} : Consacrer 10 à 15 minutes de chaque cours à la lecture silencieuse, où l'élève choisit son livre.\item \textbf{Rôle de l'enseignant} : Il ne doit pas tester la lecture (pas de quiz), mais agir comme un modèle (lire lui-même pendant ce temps) et comme un guide pour orienter les élèves vers des livres adaptés à leur niveau (règle des 98\%).\cite{III-3-2-ref44}\end{itemize}

\subsection{La stratégie de l'échelonnage (laddering)}

L'objectif final reste la lecture des textes classiques. La transition doit se faire par \og
échelons\fg{} pour éviter le découragement.

1. \textbf{Niveau Débutant (0-500 mots)} : Utilisation exclusive de \textit{Alexandros}, des
novellas très simples (\textit{O Kataskopos}), et des premiers chapitres d' \textit{Athenaze}.
L'objectif est l'automatisation des structures de base et du vocabulaire fréquent (le \og top
300\fg{}).

2. \textbf{Niveau Intermédiaire (500-2000 mots)} : Lecture des volumes suivants d'\textit{Athenaze},
des \textit{Ephodion}, et des traductions modernes de Coderch. C'est ici que l'étudiant consolide sa
morphologie verbale.

3. \textbf{Le Pont (Bridge)} : Utilisation des ressources de \textbf{Geoffrey Steadman} ou des
\textbf{Tiered Readers} de GlossaHouse. Ces ouvrages présentent des textes authentiques (Platon,
Évangiles) mais avec un vocabulaire et une grammaire expliqués sur la page en face. Cela permet de
maintenir une fluidité artificielle tout en exposant l'étudiant à la complexité
réelle.\cite{III-3-2-ref28}

4. \textbf{Lecture Authentique} : L'étudiant, ayant acquis un vocabulaire passif de 3000-4000 mots
et automatisé le parsing morphologique, peut aborder les textes originaux avec un recours minimal au
dictionnaire.



\subsection{Conclusion}

La recherche approfondie sur la Lecture Extensive appliquée au grec ancien permet de conclure sans
équivoque que cette approche n'est pas une simple option pédagogique, mais une nécessité cognitive
pour quiconque vise la fluidité.

Les preuves convergent de toutes parts :

1. \textbf{Neurosciences} : Le modèle Déclaratif/Procédural d'Ullman montre que seule la pratique
massive (input) peut transférer la compétence grammaticale vers les zones du cerveau (ganglions de
la base) capables de traiter le langage en temps réel.

2. \textbf{Psycholinguistique} : Les modèles de LaBerge \& Samuels et la théorie de la charge
fovéale expliquent pourquoi le décodage conscient (méthode traditionnelle) sature l'attention et
empêche la compréhension profonde. L'automatisation de la décomposition morphologique est la clé, et
elle ne s'acquiert que par la répétition massive offerte par les \textit{Graded Readers}.

3. \textbf{Statistique Lexicale} : Les travaux de Nation et Macdonald (Paradoxe 98-56) démontrent
mathématiquement l'impossibilité de passer directement des manuels de débutants aux textes
authentiques sans une phase intermédiaire massive de lecture graduée pour combler le déficit
lexical.

4. \textbf{Pédagogie Pratique} : L'émergence récente d'un corpus riche et varié (GlossaHouse,
novellas, méthodes inductives) a levé le dernier obstacle matériel à cette approche. Il est
désormais possible de lire des centaines de milliers de mots de grec \og facile\fg{} avant
d'affronter Thucydide.

Pour automatiser la reconnaissance sans effort, l'étudiant ne doit pas lire plus \textit{dur}, il
doit lire plus \textit{facile}, mais en quantités massives. La Lecture Extensive est le pont
technologique et cognitif qui relie enfin la connaissance théorique de la langue à sa maîtrise
vivante.

